% ECONOMICS_PROBLEM: {"id": "C73-27", "title": "Large-population selection of finite-state potential minimizers", "jel_code": "C73", "source_id": "OP-0214"}
% FIRST_POSTED: 2026-10-10
% ATTEMPT_STATUS: solved
\documentclass[11pt,a4paper]{article}
\usepackage[left=23mm,right=23mm,top=20mm,bottom=20mm]{geometry}
\usepackage[T1]{fontenc}
\usepackage[utf8]{inputenc}
\usepackage{lmodern,microtype,amsmath,amssymb,amsthm,mathtools}
\usepackage[hidelinks]{hyperref}
\usepackage{xurl}
\newtheorem{theorem}{Theorem}
\newtheorem{lemma}{Lemma}
\newtheorem{proposition}{Proposition}
\newcommand{\E}{\mathbb E}
\newcommand{\Prb}{\mathbb P}
\newcommand{\R}{\mathbb R}
\title{C73-27: exact finite-player selection\newline of a nonminimizing potential equilibrium}
\author{Ji Ho Bae}
\date{10 October 2026}
\begin{document}
\maketitle
\section*{Model, process, and claim status}
\noindent\textbf{Model:} GPT Codex. \textbf{Contributor:} Ji Ho Bae.
The exact serving revision was not exposed.
The complete mathematical argument below was generated by GPT Codex,
using parallel derivation and adversarial proof checking in response to a
request for a complete hand proof of the exact catalogue statement.
No human-written mathematical edits are included. The construction and
all verification arguments are preserved in full.
This is an LLM claim of a complete counterexample; community expert review
is pending. No proof-assistant verification is claimed.
The complete original generated output is preserved at the immutable
\href{https://github.com/jbaelaw/finite-state-nonminimizing-selection/blob/58b261b80171ab6a1ddb94a9dab52ab6b50c0c3d/paper/Proof.tex}{manuscript TeX}
and \href{https://github.com/jbaelaw/finite-state-nonminimizing-selection/blob/58b261b80171ab6a1ddb94a9dab52ab6b50c0c3d/paper/Proof.pdf}{manuscript PDF}.

\begin{abstract}
Fixed smooth three-state data give bounded symmetric exact Markov-perfect
feedbacks for every population size. Their empirical paths converge to
a constant equilibrium that does not minimize the potential action.
We verify every continuation configuration and all bounded progressively
measurable unilateral deviations, prove the full path limit, and exhibit
a strictly better feasible path.
\end{abstract}

\section{Statement and smooth data}
Write $\Delta_3=\{m\in[0,1]^3:\sum_i m_i=1\}$ and fix
$T=1$, $m_0=(1/3,1/3,1/3)$. For the potential action in C73-27,
\begin{equation}\label{action}
 J(m,a)=\int_0^1\left[\frac12\sum_{i,j\ne i}m_i a_{ij}^2+
 \mathcal F(m)\right]dt+\mathcal G(m(1)),\qquad
 \dot m_i=\sum_{j\ne i}(m_ja_{ji}-m_ia_{ij}),
\end{equation}
the admissible paths are absolutely continuous simplex paths starting
at $m_0$ with measurable nonnegative off-diagonal rates and finite
weighted quadratic energy. Denote the paths in global minimizers by
$\mathcal M_{1,m_0}$. In the finite game, player $\ell$'s terminal cost is
$g_{X_1^\ell}(m_1^{N,-\ell})$, where $g_i=\partial_i\mathcal G$ and
$m^{N,-\ell}$ is the empirical law of the other $N-1$ players.
The running cost is $\frac12\sum_{j\ne X_t^\ell}(a^\ell_{X_t^\ell j})^2+
f_{X_t^\ell}(m_t^{N,-\ell})$, with $f_i=\partial_i\mathcal F$.

\begin{theorem}\label{main}
There are $C^\infty$ potentials on $\R^3$ and, for every $N\ge2$,
bounded symmetric Markov-perfect feedbacks for exactly these games
whose empirical processes satisfy
\[
 \sup_{0\le t\le1}\|\mu_t^N-m_0\|_2\longrightarrow0
 \quad\text{in probability},\qquad
 \mathcal L(\mu^N)\Longrightarrow Q=\delta_{\bar m}
 \quad\text{in }D([0,1],\Delta_3)\text{ with }J_1,
\]
where $\bar m(t)=m_0$. This continuous limit has
$Q(\mathcal M_{1,m_0})=0$. In particular, the universal selection
assertion in the stated C73-27 formulation is false.
\end{theorem}

Set $a=2/5$, $b=3/5$ and define
\[
 \rho(r)=\begin{cases}e^{-1/r},&r>0,\\0,&r\le0,\end{cases}
 \qquad
 \psi(z)=\frac{\rho(z-a)}{\rho(z-a)+\rho(b-z)}.
\]
The denominator is positive everywhere. Every right derivative of
$e^{-1/r}$ is a polynomial in $1/r$ times $e^{-1/r}$ and tends to zero
at $0$, so $\rho$ and $\psi$ are smooth. The quotient rule gives
\[
 \psi'(z)=\frac{\rho'(z-a)\rho(b-z)+\rho(z-a)\rho'(b-z)}
 {[\rho(z-a)+\rho(b-z)]^2}\ge0.
\]
Thus $\psi=0$ on $(-\infty,a]$, $\psi=1$ on $[b,\infty)$,
and $L:=\sup_\R\psi'<\infty$. With $S(m)=\sum_i m_i$, choose
\begin{equation}\label{data}
 \mathcal F(m)=0,\qquad
 \mathcal G(m)=-\psi(m_1)+(S(m)-1)\psi'(m_1).
\end{equation}
These are fixed independently of $N$. Direct ambient differentiation
yields $\partial_1\mathcal G=(S-1)\psi''(m_1)$ and
$\partial_2\mathcal G=\partial_3\mathcal G=\psi'(m_1)$. Consequently
on the simplex the actual data are
\begin{equation}\label{costs}
 \mathcal G(m)=-\psi(m_1),\qquad
 (g_1,g_2,g_3)=(0,\psi'(m_1),\psi'(m_1)),\qquad f_i=0.
\end{equation}
All individual terminal costs lie in $[0,L]$.

\section{An exact feedback for every finite game}
For $x\in\{1,2,3\}^N$ let
$k_\ell(x)=\#\{r\ne\ell:x_r=1\}$ and write $x^{\ell\to j}$ for
the vector with coordinate $\ell$ replaced by $j$.
Construct scalar functions in reverse time $\tau\in[0,1]$ by
\begin{align}
 w_{N-1}'&=-\tfrac12w_{N-1}^2,& w_{N-1}(0)&=\psi'(1)=0,\label{ode-top}\\
 w_k'&=(N-1-k)w_k w_{k+1}-(N-k-\tfrac12)w_k^2,
 &w_k(0)&=\psi'\bigl(k/(N-1)\bigr),\quad 0\le k<N-1.
 \label{ode}
\end{align}
\begin{lemma}\label{ode-lemma}
This system has a unique $C^1$ solution with $0\le w_k\le L$.
If $\psi'(k/(N-1))=0$, then $w_k\equiv0$.
\end{lemma}
\begin{proof}
The top function is zero. Proceed downwards, assuming
$0\le w_{k+1}\le L$. Put $A_k=(N-1-k)w_{k+1}$,
$B_k=N-k-1/2$ and $c_k=w_k(0)$. The equation
$w_k'=A_kw_k-B_kw_k^2$ is locally Lipschitz in $w_k$.
For $c_k=0$ its unique solution is zero. For $c_k>0$ its global solution is
\[
 w_k(\tau)=\frac{c_k\exp(\int_0^\tau A_k(r)dr)}
 {1+B_kc_k\int_0^\tau\exp(\int_0^s A_k(r)dr)ds}.
\]
Differentiation verifies the equation and initial value. The denominator
is positive. Whenever $w_k>L$, its derivative is bounded above by
$(N-1-k)Lw_k-(N-k-1/2)w_k^2<0$. A connected interval above $L$
would therefore be strictly decreasing from its entry value $L$,
a contradiction. This proves the bound and the induction.
\end{proof}

Define values and outgoing feedbacks by
\begin{equation}\label{values}
 \begin{aligned}
 v_\ell(t,x)=\begin{cases}0,&x_\ell=1,\\
 w_{k_\ell(x)}(1-t),&x_\ell\in\{2,3\},\end{cases}
 \\[3pt]
 q_j^\ell(t,x)=\begin{cases}
 w_{k_\ell(x)}(1-t),&x_\ell\in\{2,3\},\ j=1,\\
 0,&\text{otherwise},
 \end{cases}\quad j\ne x_\ell.
 \end{aligned}
\end{equation}
The diagonal is minus the outgoing sum. Rows for counterfactual own
states are given by the same rule after replacing $x_\ell$ by that
state, so these define admissible full rate matrices. They are continuous
in time, bounded by $L$ independently of $N$, and symmetric under
player permutations. Independent Poisson transition noises implement
the joint chain; its total intensity is at most $NL$, so it is
nonexplosive and uses no common noise.

The exact Nash equations are
\begin{equation}\label{hjb}
 \partial_t v_\ell+
 \sum_{r\ne\ell}\sum_{j\ne x_r}q_j^r
 [v_\ell(t,x^{r\to j})-v_\ell(t,x)]
 -\tfrac12\sum_{j\ne x_\ell}
 [v_\ell(t,x)-v_\ell(t,x^{\ell\to j})]_+^2=0.
\end{equation}
Their terminal condition is exactly
$v_\ell(1,x)=g_{x_\ell}(m^{N,-\ell}_x)$.
To verify \eqref{hjb} at every configuration, first suppose $x_\ell=1$.
Its value is zero under all opponent changes, and own alternative
values are nonnegative. Thus all terms vanish and its optimal rates
are zero. If $x_\ell\ne1$, set $k=k_\ell(x)$ and $\tau=1-t$.
Its own jump to $1$ changes its value from $w_k$ to $0$; switching
between $2$ and $3$ leaves its value unchanged. Its own Hamiltonian
is therefore $w_k^2/2$. Of the opponents, the $k$ in state $1$ have
zero rates. Each of the $N-1-k$ others has exactly $k$ state-$1$
opponents and jumps to $1$ at rate $w_k$, changing the focal value
to $w_{k+1}$. Equation \eqref{hjb} becomes
$-w_k'+(N-1-k)w_k(w_{k+1}-w_k)-w_k^2/2=0$, precisely
\eqref{ode}; when $k=N-1$ it is \eqref{ode-top}. In all cases,
\begin{equation}\label{optimal-rate}
 q_j^\ell(t,x)=[v_\ell(t,x)-v_\ell(t,x^{\ell\to j})]_+.
\end{equation}

\begin{proposition}\label{mpe}
The feedbacks \eqref{values} are exact Markov-perfect Nash equilibria
against every bounded progressively measurable unilateral rate deviation.
\end{proposition}
\begin{proof}
Fix any initial continuation time $t_0$, state vector $x$ and player
$\ell$, retaining the opponents' feedbacks. Let its outgoing rates
$\alpha_j(t)\ge0$ be any admissible bounded progressively measurable
deviation, possibly depending on the entire history, and let $X$ be
the resulting state process. Put
$\Delta_j=v_\ell(t,X_t^{\ell\to j})-v_\ell(t,X_t)$.
The compensated-jump formula has a mean-zero martingale because the
values, their time derivatives for fixed $N$, and all intensities
are bounded. Using \eqref{hjb} in that formula gives exactly
\begin{align}\label{verify}
 &\E\left[\int_{t_0}^1\tfrac12\sum_{j\ne X_t^\ell}\alpha_j(t)^2dt
 +g_{X_1^\ell}(m_1^{N,-\ell})\right]-v_\ell(t_0,x)\nonumber\\
 &\quad=\E\int_{t_0}^1\sum_{j\ne X_t^\ell}
 \left[\tfrac12(\alpha_j-q_j^\ell)^2+\alpha_j(\Delta_j)_+\right]dt\ge0.
\end{align}
Indeed, for $\alpha\ge0$ and $q=(-\Delta)_+$,
$\alpha^2/2+\alpha\Delta+q^2/2
=(\alpha-q)^2/2+\alpha\Delta_+$.
By \eqref{optimal-rate} the prescribed feedback makes this remainder
zero. It therefore attains the minimum against all allowed deviations.
The choices of $t_0,x,\ell$ were arbitrary, proving the continuation
condition at every time and configuration, including rare states.
\end{proof}

\section{Convergence of the entire empirical path}
Take the required independent initial states with law $m_0$, independent
of the transition noises. Let $K_N$ count those initially in state $1$,
and set $A_N=\{K_N\le(2/5)(N-1)\}$. On $A_N$, every player outside
state $1$ has $K_N$ state-$1$ opponents, so its initial scalar terminal
datum is zero. Lemma~\ref{ode-lemma} gives
$w_{K_N}(1-t)=0$ at every time. Players in state $1$ also have zero
rates. This configuration is absorbing under the exact feedback,
and $\mu_t^N=\mu_0^N$ throughout $[0,1]$ on $A_N$.

Here $K_N$ is binomial with mean $N/3$ and variance $2N/9$.
For $N\ge12$ the threshold exceeds its mean by at least $N/30$;
Chebyshev's inequality gives $\Prb(A_N^c)\le200/N$. Also
$\E\|\mu_0^N-m_0\|_2^2=2/(3N)$.
Since the squared simplex diameter is at most $2$,
\begin{equation}\label{convergence}
 \E\sup_{t\le1}\|\mu_t^N-m_0\|_2^2
 \le\frac{2}{3N}+2\Prb(A_N^c)
 \le\frac{2}{3N}+\frac{400}{N}\longrightarrow0.
\end{equation}
Markov's inequality proves uniform convergence in probability.
Choosing the identity time change bounds the $J_1$ distance by the
uniform distance, so the laws converge in precisely the topology in
the question to $Q=\delta_{\bar m}$. The limiting path is continuous.

\section{Existence of minimizers and strict nonminimality}
For completeness, the minimizing-path set is nonempty and Borel.
Introduce fluxes $z_{ij}=m_i a_{ij}\ge0$. The constraints become linear
in $z$, and the energy is
$K(m,z)=\frac12\int_0^1\sum z_{ij}^2/m_i\,dt$, with $0/0=0$ and
$z^2/0=+\infty$ for $z>0$. The constant zero-flux path is feasible.
Since the terminal potential belongs to $[-1,0]$, a minimizing sequence
has bounded energy, hence bounded $L^2$ fluxes because $m_i\le1$.
The forward equation gives bounded $L^2$ path derivatives and a common
$1/2$-H\"older modulus. Arzel\`a--Ascoli and weak $L^2$ compactness
give uniform $m^n\to m$ and weak $z^n\rightharpoonup z$ along a
subsequence. Nonnegative test functions preserve $z\ge0$, and time
interval indicators pass the integrated constraints to the limit.

The exact dual representation
\[
 K(m,z)=\sup_{q\in L^\infty}\int_0^1\sum_{i,j\ne i}
 [q_{ij}z_{ij}-\tfrac12m_iq_{ij}^2]dt
\]
proves lower semicontinuity: each bounded $q$ gives a continuous
expression under these convergences. Equality in the representation
follows by the bounded choices $q_{ij}=\min\{A,z_{ij}/m_i\}$,
assigning $A$ where $m_i=0<z_{ij}$ and zero where both vanish, and
monotone convergence as $A\uparrow\infty$. Continuity passes the
terminal potential to the limit. Thus the limit minimizes. Finite
energy forces $z_{ij}=0$ almost everywhere on $\{m_i=0\}$; defining
$a_{ij}=z_{ij}/m_i$ on $\{m_i>0\}$ and zero elsewhere recovers
admissible measurable rates. The same argument for any sequence of
minimizers shows that their law paths form a uniformly compact set.
Its inclusion in $D$ is continuous for $J_1$, so it is compact and
Borel there as well.

Every rate representation of $\bar m$ has nonnegative action, since
$\mathcal F=0$, $\mathcal G(m_0)=0$, and kinetic energy is nonnegative.
For a competitor choose $c=\log2$, $a_{21}=a_{31}=c$, and all other
off-diagonal rates zero. The forward law is
\[
 m_1(t)=1-\tfrac23e^{-ct},\qquad
 m_2(t)=m_3(t)=\tfrac13e^{-ct}.
\]
It is admissible and has terminal law $(2/3,1/6,1/6)$, where
$\mathcal G=-1$. Its energy is
$\frac12\int_0^1(m_2+m_3)c^2dt=c(1-e^{-c})/3=(\log2)/6$.
Therefore
\begin{equation}\label{strict}
 J(m,a)=\frac{\log2}{6}-1<0.
\end{equation}
For example $\log2=\int_1^2x^{-1}dx<1$ proves this strict inequality
without computation. The global minimum is negative, so $\bar m$
cannot occur in any minimizing pair. This gives
$Q(\mathcal M_{1,m_0})=0$ and completes Theorem~\ref{main}.

\section{The prescribed representatives and source scope}
The normal term in \eqref{data} vanishes on the simplex but adds the
common component $\psi'(m_1)$ to the ambient gradient of $-\psi(m_1)$.
Its tangent projection is $(-2\psi'/3,\psi'/3,\psi'/3)$, exactly the
intrinsic gradient of the restricted potential. This is allowed by
the ambient derivative prescription in C73-27 and by the projected
potential convention in \cite[Eq.~(2.9)]{CeD}; no hypothesis requires
$\sum_i g_i=0$. We retain the actual costs \eqref{costs} in every
finite continuation game. State-dependent common additions are not
asserted to preserve finite closed-loop feedbacks. Centering those
costs would change the finite game and would require another analysis.

The published two-state selection result cited in the background
analyzes a specified anti-monotone model, rather than all smooth
terminal costs: \cite[Section~3]{CDFP} uses $G(x,m)=-mx$ on
$\{-1,1\}$ with zero interaction running cost. The potential
selection discussion in \cite[Section~4.3]{ICM} supplies the motivation;
the theorem here addresses the exact, explicitly quantified catalogue
formulation. It uses all required finite-player continuation incentives,
the stipulated independent noises and initial law, the full path
limit, and the original global potential action.

\begin{thebibliography}{9}
\bibitem{ICM} P. Cardaliaguet and F. Delarue,
\emph{Selected topics in mean field games}, Proc. ICM 2022, Vol.~5,
EMS Press (2023), 3660--3703, doi:10.4171/ICM2022/17,
Section~4.3, equations (4.18)--(4.19).
\url{https://ems.press/content/book-chapter-files/33265}.
\bibitem{CeD} A. Cecchin and F. Delarue,
\emph{Selection by vanishing common noise for potential finite state
mean field games}, Comm. Partial Differential Equations 47 (2022),
89--168, doi:10.1080/03605302.2021.1955256, equation (2.9).
\url{https://arxiv.org/abs/2005.12153}.
\bibitem{CDFP} A. Cecchin, P. Dai Pra, M. Fischer and G. Pelino,
\emph{On the convergence problem in mean field games: a two state
model without uniqueness}, arXiv:1810.05492v2 (2019), Section~3.
\url{https://arxiv.org/abs/1810.05492}.
\end{thebibliography}
\end{document}
